\chapter{Conclusion and Outlook}
\label{ch:conclusion}

\section{Conclusion}
\label{sec:conclusion}

This thesis set out to identify plasma biomarkers of subclinical atherosclerosis and ended by
examining the phenotype such a search depends on. Three findings summarise the work.

\paragraph{The pipeline was reproduced faithfully, and the phenotype it yields carries little
detectable proteomic signal.}
Re-running the detection pipeline of \textcite{omarov2024deep} across the UK Biobank carotid
ultrasound collection produced detections for 179{,}611 images from 20{,}014 participants and a
plaque prevalence of 43.0\,\% at participant level, against the 45\,\% reported by the original
authors. Tested against the Olink proteomics panel in the 2{,}542 participants for whom both were
available, this phenotype showed no association that survived correction for multiple testing:
151 of 2{,}923 proteins reached nominal significance where approximately 146 are expected by
chance, and the best-performing classifier reached a cross-validated ROC-AUC of 0.541, below the
0.551 attained by age and sex alone. Age itself was significantly associated after correction and
produced a visible gradient in the leading principal components, confirming that the analysis
detects association where association exists.

\paragraph{The base model transfers to an external dataset with a substantial loss of
specificity, and fine-tuning recovers part of it.}
On the 48 held-out images of Ar-PlaqSegm1 \cite{rulloni2025arplaqsegm1}, the base model
reached a sensitivity of 83.3\,\% and a specificity of 62.5\,\%, against 89.5\,\% and 89.2\,\%
on its own UK Biobank test set. Sensitivity therefore declined by 6.2 percentage points and
specificity by 26.7. Four segmentation models fine-tuned on the Ar-PlaqSegm1 development set
were individually indistinguishable from one another on the test set, each reaching a Youden's
$J$ of 0.625 against the base model's 0.458 and differing only in how they distributed their
errors. Pooling the predictions of two of them raised $J$ to 0.667 at a sensitivity of 91.7\,\%
and a specificity of 75.0\,\%, recovering roughly half of the specificity lost to the change of
data source. All operating points were selected on the validation split rather than on the test
set; had they been optimised on the test set, the same configuration would have reported 0.792.

\paragraph{The limiting component is confidence ranking, not segmentation.}
Taking the best-matching predicted instance for each plaque-positive test image, irrespective of
its confidence score, raises the mean Dice coefficient from 0.655 to 0.794 and leaves no image
without a correct mask. Both apparent failures of the pooled model disappear: the most severe is
outlined at a Dice of 0.82 while carrying a confidence of 0.0019, against an operating threshold
of 0.25. The models locate plaques they do not report. This behaviour was traced to the weighting
of the training objective, in which the classification term -- the confidence score, for a single
class -- was weighted twenty-five times lower than the terms governing geometry. Raising that
weight lifted the median confidence of the correct instance from 0.060 to 0.602 but degraded mask
quality, placing the optimum between the two settings tested. Attempting to repair the ranking
after training improved the discrimination of individual instances substantially, from an area
under the curve of 0.681 to 0.823, yet lowered image-level performance: the features that
identify a genuine plaque also promote plaque-shaped predictions in plaque-free images. The
binding constraint is therefore objectness rather than instance ranking, and it is not
addressable post hoc.

\section{Outlook}
\label{sec:conclusion-outlook}

\paragraph{Continuing the biomarker search.}
The analyses of Phases 1A and 1B remain open and can be resumed once access to the UK Biobank
Research Analysis Platform is restored, or if an alternative cohort linking carotid imaging to
proteomics becomes available. The objective that motivated the original design is the one still
worth pursuing: a health-record-based proxy for the imaging phenotype, fitted where imaging and
records coexist and then evaluated wherever records exist. It would lift the proteomic analysis
from 2{,}542 participants to the full Olink set, and a null result at that sample size would
carry weight that a null result at 2{,}542 cannot.

Two considerations govern whether such a proxy would work. The first is its own accuracy: the
proxy inherits the error of the imaging phenotype it is fitted against and adds its own, so the
attenuation argument of Section~\ref{sec:domain-shift} applies twice over, and the ceiling on
what it can detect is set by the model examined in the third part of this thesis. The second is
specificity. A proxy built from age, blood pressure, lipids and diagnostic codes predicts plaque
through the risk factors that cause it, and plasma proteins are themselves associated with those
risk factors; an association between proxy and proteome could therefore reflect the risk-factor
profile rather than the lesion. Establishing that the proxy carries information about plaque
beyond its inputs -- for instance by testing whether it improves on those inputs alone within the
2{,}542 participants where the imaging phenotype is available -- is a prerequisite for
interpreting anything found at the larger sample size.

Two smaller changes to the original design are also warranted. The clinical baseline should use
age at the imaging visit rather than age at recruitment, since the plaque phenotype derives from
the former and the interval between the two varies between participants. And the comparison
against an independent risk estimate, which Phase 1B was also to provide, would still help
determine whether the null result originates in the phenotype or in the proteomic data.

\paragraph{Avoiding the discrete phenotype altogether.}
The attenuation argument that motivates this thesis applies to a pipeline that first reduces an
image to a binary label and then correlates that label with molecular data. End-to-end multimodal
fusion, in which image and tabular data enter a single model jointly, never forms the discrete
phenotype and therefore does not incur the associated loss of information. The approach is
established across medical imaging and health-record data more generally
\cite{huang2020fusion}, and \textcite{stamos2025} reports an instance of it for carotid
plaque risk stratification, with imaging contributing the greater part of the signal and tabular
data a smaller complement. Applying it at cohort scale, with the imaging arm initialised from a
model of the kind developed here, is a natural continuation.

\paragraph{Making results comparable across studies.}
Section~\ref{sec:discussion-literature} sets out why the published figures for carotid plaque
detection cannot be compared with one another: each was obtained on data held privately by the
group that produced it. The evaluation reported here differs in that its test partition is
public, so the figures can be recomputed and any later model placed beside them without
renegotiating data access. Generalising that is less a methodological problem than a matter of
convention. Several openly available carotid ultrasound datasets, from different centres, devices
and case mixes, together with the practice of reporting on all of them, would turn the aggregate
domain-shift estimate obtained here into a decomposition by factor: how much of the loss is the
device, how much the population, how much the operator. Evaluating on one external dataset
establishes that transfer is imperfect; evaluating on several would establish why.

\paragraph{Establishing which differences are real.}
The four fine-tuned models could not be separated on 48 test images, in which one image
corresponds to 4.2 percentage points of sensitivity. Conclusions about the classification loss
weight, the choice of architecture or the value of ensembling therefore rest on differences that
the available evaluation set cannot resolve. Cross-validation over the full development set would
provide a mean and a spread instead of a single point estimate, at a cost of roughly five
training runs.

\paragraph{Encoding the diagnostic criterion.}
The false positives that both models share follow the vessel wall over an extended stretch
rather than enclosing a focal protrusion, and are systematically thinner than true plaque
annotations. This is consistent with a model that classifies each candidate region on its own
appearance, whereas the diagnostic criterion is relative: a plaque exceeds 50\,\% of the
\emph{surrounding} intima-media thickness. Supplying that comparison explicitly -- by segmenting
the wall as a second class, or by providing local wall thickness as an additional input -- would
target the dominant failure mode directly, rather than hoping the model infers the criterion from
examples.

\paragraph{Using the discarded frames.}
Main-axis selection reduced 179{,}611 UK Biobank images to 42{,}954, discarding 76\,\% of the
material. Carotid ultrasound acquisitions are cine loops, and the discarded frames carry motion
information that a single-frame model cannot use. Spatiotemporal approaches to plaque
characterisation have been proposed but are computationally demanding; whether they repay that
cost at cohort scale is untested.

\paragraph{Extending the annotated material.}
The one test image that no model segments correctly even under ideal confidence ranking, and the
smallest plaque in the test set, both point to under-represented regions of the training
distribution: plaques in the smallest 6\,\% of the annotated size range are supported by only 28
training instances. Targeted annotation of small and atypical plaques would address this more
efficiently than uniformly enlarging the dataset.

\paragraph{Closing remark.}
This thesis set out to look for molecular markers of a disease and ended by examining how that
disease is measured. The two are not separate undertakings. An association study is bounded by
the quality of the variables it relates, and where one of those variables is produced by a model
rather than observed, the error of that model becomes part of the epidemiology. The proteomic
analysis returned no signal, and nothing established here settles whether that is because no
signal exists or because the phenotype was too coarse to carry one. What the work does establish
is how large the measurement gap becomes outside the cohort the model was built in, which part of
it is the representation and which the threshold, and that in one specific respect the gap is
smaller than it appears: the models outline plaques they decline to report, and the deficit lies
in how the training objective was weighted rather than in what the network can see. Examining the
instrument before trusting what it measures is not a detour from the biological question; it is
the step the answer rests on.
